\textbf{COLUNI 2018 -- Questão 37}\par\medskip Leia o texto abaixo:

\par

“Na década de 1960, intelectuais, escritores, cantores, compositores, atores, cineastas e artistas plásticos participaram ativamente da política e muitos se engajaram em projetos revolucionários. A renovação artística tornou-se muito politizada após o golpe militar. Com o movimento sindical reprimido, as esquerdas investiram na produção cultural.” (VAINFAS, Ronaldo. et al. \textbf{História}. v.3. 2.ed, São Paulo: Saraiva, 2013, p. 222)

\par

Sobre o contexto cultural da década de 1960 no Brasil, é INCORRETO afirmar que:\begin{enumerate}
\item a União Nacional dos Estudantes (UNE) criou os Centros Populares de Cultura (CPCs) que produziram
        música e teatro para conscientizar a população a respeito da realidade econômica e social do país.
\item Caetano Veloso, Gilberto Gil e Torquato Neto criaram o movimento musical conhecido como Tropicália que
        juntava à música brasileira elementos do movimento da contracultura que ocorria nos EUA como o rock.
\item o Compositor Chico Buarque de Holanda produziu obras teatrais como “Calabar, o elogio da Traição” e
        várias músicas que foram censuradas pelo Regime Militar por criticar a história tradicional do Brasil.
\item o movimento musical da “Jovem Guarda” com a participação de Roberto Carlos, Erasmo Carlos e
        Wanderléia incitava os jovens à mobilização política e era o maior contestador do Regime Militar.\end{enumerate}
